\weeknum{9}
\topic[Transient Heat Flow]{Geothermics: Transient Heat Flow}

\workshopstart

\vspace{-0.75em}
\section*{Preparation}

\begin{description}
  \item[Pre-reading:] Course Notes, Ch. 11
  \item[Lecture videos:]
  \item[Notes:]
\end{description}

\section*{Purpose}

This workshop develops intuition for transient heat flow: how thermal diffusivity sets the timescales and length scales over which thermal signals propagate, how half-space and plate cooling models describe the thermal evolution of oceanic lithosphere, and how latent heat modifies cooling in Stefan problems.

\section*{Learning Outcomes}
Students will:
\begin{itemize}
  \item Understand how thermal diffusivity controls timescales.
  \item Compare thermal and chemical diffusion processes.
  \item Interpret half-space and plate cooling models.
  \item Use Fourier series to conceptualize thermal evolution.
\end{itemize}

\section*{Session Flow (\timelinelength\ min)}

% {\color{red}
% \timeline[slot]{Warm-up --- How long does it take to cook a roast?}{10}

% Ask students: \emph{Why does a larger roast take longer to cook in the oven even at the same temperature? What does this tell us about heat moving through rock?} Elicit the idea that the time for heat to penetrate scales with the \emph{square} of the distance. Connect to the diffusion equation and introduce thermal diffusivity $\kappa$ as the quantity that controls this scaling.
% }

\timeline[warmup]{Wedge vs.\ Orogen}{10}

Both accretionary wedges and orogens form at convergent margins, but they experience very different thermal histories. Before we build a detailed model for one of them later today, predict how surface heat flow evolves over the life of each.

\timeline[slot]{Mini-lecture --- The diffusion equation and thermal diffusivity}{15}

\begin{itemize}
  \item The heat (diffusion) equation: $\partial T/\partial t = \kappa\, \nabla^2 T$.
  \item Thermal diffusivity $\kappa = k/(\rho c_p)$: what controls it, typical values for crustal rocks.
  \item Scaling argument: thermal length $z \sim \sqrt{\kappa t}$; thermal timescale $t \sim z^2/\kappa$.
  \item Key difference from steady state: temperature \emph{evolves} in time.
\end{itemize}

\timeline[item]{Transient Geotherms}{40}

Tectonic and other geologic events often have a thermal effect on the lithosphere, whether top, bottom, or internally.  In this activity, you'll explore several scenarios related and learn how to approximate their temporal evolution intuitively.


\timeline[item]{Thermal Diffusivity, Thermal Length, and Timescales}{45}
\begin{itemize}
  \item Review $\kappa = k / (\rho c_p)$.
  \item Emphasize modest variation in crustal rocks.
  \item Contrast sediments, rock, and ice.
  \item Penetration depth scaling: $z \sim \sqrt{\kappa t}$.
  \item Borehole paleothermometry vs ice-core records.
  \item Depth--resolution tradeoffs.
\end{itemize}

% \timeline[item]{Thermal vs Chemical Diffusion}{30}
% \begin{itemize}
%   \item Both obey $\sqrt{Dt}$ scaling.
%   \item Chemical diffusion often multi-component.
%   \item Zoning persistence vs thermal smoothing.
% \end{itemize}

\timeline[slot]{Break}{10}

\timeline[item]{Oceanic Lithosphere: Half-Space and Plate Cooling}{35}
\begin{itemize}
  \item Derive surface heat flow and seafloor subsidence from the half-space cooling solution.
  \item Test both predictions against real global heat flow and bathymetry data.
  \item Compare to the plate cooling model to see where and why the half-space prediction breaks down, at $t\sim\tau=L^2/\kappa$.
\end{itemize}

% \timeline[item]{Stefan Problems in Geology}{30}
% \begin{itemize}
%   \item Latent heat as effective heat capacity.
%   \item Cooling times of dikes vs sills.
%   \item When analytic solutions break down.
% \end{itemize}

\timeline[slot]{Wrap-up}{10}

\emph{What single quantity controls how fast heat moves through rock, and what are its typical values?}
Nearly all thermal problems result in a time or distance scaling that comes from the argument of an error function or exponential similar to $t \sim z^2/\kappa$.  Can you give one geological example from each activity where this scaling is evident?

\timeline[slot]{Preview: Finite Differences}{10}

Analytic solutions to the diffusion equation require simplifying assumptions: constant properties, simple geometries, and idealized boundary conditions. Real geological systems violate all of these Next week we develop finite difference methods, which replace the differential equation with a local energy balance that can be applied to arbitrarily complex geometries and boundary conditions. You will first derive the finite difference equations from first principles and then solve the heat equation yourself by hand---becoming the computer.

\workshopend
\notepage
\notepage

\subimport{activities/}{warmup_collision}
\subimport{activities/}{activity_transient_scenarios}
\subimport{activities/}{activity_thermal_timescales}
\subimport{activities/}{activity_oceanic_cooling}